\begin{abstract}
Register tokens were originally introduced to remove high-norm attention artifacts in large vision transformers trained on massive datasets. Their behavior in data-scarce regimes, where a compact model is fine-tuned on limited samples, remains largely uncharacterized. In this work, we investigate whether register tokens provide structural regularization under data scarcity. We conduct a fully paired, controlled empirical ablation across 24 training runs on CIFAR-100 subsets, evaluating ViT-Tiny across four register allocations ($K \in \{0, 1, 4, 8\}$) and two data budget regimes (100 and 300 images per class, with three random seeds per arm). On the untouched 10{,}000-image test set at 100 images per class, no register configuration outperforms the register-free control ($75.19 \pm 0.24\%$). While $K=4$ exhibits an apparent regularization signal on the small validation split by lowering validation loss ($\Delta = -0.034$, $p = 0.024$) and narrowing the generalization gap, this effect fails to transfer to held-out test performance ($\Delta = +0.002$, $p = 0.823$). To test whether this null result stems from extreme data starvation, we scale the training budget threefold to 300 images per class. While baseline test accuracy surges by $+6.14\%$ to $81.33 \pm 0.37\%$, register tokens continue to yield no statistically significant advantage ($81.23 \pm 0.19\%$ for $K=1$, $81.23 \pm 0.18\%$ for $K=4$, and $81.02 \pm 0.45\%$ for $K=8$). Furthermore, layer-wise attention entropy and patch-norm outlier rates confirm the absence of attention-sink suppression. We conclude that register tokens do not regularize vision transformers under data scarcity, and that mechanisms observed at foundation scale do not automatically generalize to compact models in low-data regimes.
\end{abstract}
